% Why the timestamps have to be reconstructed rather than assumed.
\begin{tikzpicture}[font=\sffamily\scriptsize, x=1cm, y=1cm]

% ================================================================ band 1
\node[signame] at (-0.3,2.9) {the sensor};
% samples on the sensor's own clock, evenly spaced
\foreach \i in {0,...,15}{ \node[sample] at (0.3+\i*0.62,2.9) {}; }
\draw[busline] (0.2,2.9) -- (9.7,2.9);
\draw[faultedge] (9.6,3.35) -- (9.6,3.0);
\node[tlabel, anchor=south, text=red!60!black] at (9.6,3.4) {watermark: the pin rises};

\node[signame] at (-0.3,1.9) {the batch, as read};
\node[taskrun, minimum width=58mm, anchor=south west, fill=hwcol] at (4.1,1.7) {all sixteen samples arrive at once};
\node[tlabel, anchor=west] at (10.0,1.9) {and none of them carries a time};

% ================================================================ band 2
\node[signame] at (-0.3,0.6) {reconstructed};
\foreach \i in {0,...,15}{ \draw[busstub] (0.3+\i*0.62,0.35) -- (0.3+\i*0.62,0.85); }
\node[tlabel, anchor=west] at (10.0,0.6) {newest = the capture};
\node[tlabel, anchor=west] at (10.0,0.2) {older = capture minus k intervals};

% ================================================================ band 3
\node[signame] at (-0.3,-0.8) {control periods};
\foreach \i in {0,...,9}{
  \draw[release] (0.15+\i*1.05,-1.25) -- (0.15+\i*1.05,-0.55);
  \draw[tick] (0.15+\i*1.05,-1.45) -- (0.15+\i*1.05,3.2);
}
\node[tlabel, anchor=west] at (10.0,-0.8) {on a different clock again};

\node[signame] at (-0.3,-2.0) {what the loop takes};
\foreach \i in {0,...,9}{
  \node[sample] at (0.15+\i*1.05+0.18,-2.0) {};
  \draw[tspan] (0.15+\i*1.05,-2.35) -- (0.15+\i*1.05+0.18,-2.35);
}
\node[tlabel, anchor=west] at (10.0,-2.0) {the nearest sample, and its age};

% ================================================================ band 4
\node[chlabel, anchor=north west] at (-0.3,-3.0) {and the two clocks drift apart};
\draw[taxis] (0.2,-4.6) -- (9.7,-4.6);
\draw[signal, line width=0.9pt] (0.2,-4.55) -- (9.7,-3.7);
\node[tlabel, anchor=west] at (9.8,-3.7) {error using the nominal interval};
\draw[signal, line width=0.9pt, draw=blue!55!black] (0.2,-4.55) -- (9.7,-4.45);
\node[tlabel, anchor=west] at (9.8,-4.45) {error using the measured interval};
\node[tlabel, anchor=north] at (5.0,-4.7) {time, over one hour};

\node[umlnote, text width=56mm, anchor=north west] at (-0.3,-5.3)
  {Three different clocks appear in this figure: the sensor's own oscillator,
   the node's timer, and, once chapter 12 arrives, the bus master's. The node
   converts everything into its own time as early as possible, which here means
   at the moment the batch is read, and never carries a timestamp in somebody
   else's units.};
\node[umlnote, text width=52mm, anchor=north west] at (6.4,-5.3)
  {The lower band is why step 4 exists. The nominal interval is a number from a
   datasheet; the measured one is what this sensor achieves on this bench today.
   Using the first makes the oldest sample in every batch progressively wrong,
   and the error is invisible on any plot that shows only the newest sample.};
\end{tikzpicture}
